\def\tituloconcepto{Oscilación de Rabi}
% Si el capitulo se compila solo (sin main.tex), esta cabecera abre el documento.
% Antes de % Si el capitulo se compila solo (sin main.tex), esta cabecera abre el documento.
% Antes de % Si el capitulo se compila solo (sin main.tex), esta cabecera abre el documento.
% Antes de \input{cabecera} cada capitulo define \tituloconcepto.
\ifdefined\enmain\else
  \documentclass[11pt,a4paper]{article}
  \input{preambulo}
  \begin{document}
  \title{\tituloconcepto}
  \author{\autor}
  \date{\today}
  \maketitle
\fi
 cada capitulo define \tituloconcepto.
\ifdefined\enmain\else
  \documentclass[11pt,a4paper]{article}
  % Preambulo compartido por todos los documentos. Compilar con XeLaTeX/LuaLaTeX
% (o Tectonic). Si se usa Overleaf, elegir el compilador XeLaTeX.
\usepackage[spanish,es-noshorthands]{babel}
\usepackage{amsmath,amssymb}
\usepackage{booktabs,array}
\usepackage[a4paper,margin=2.5cm]{geometry}
\usepackage{xcolor}
\usepackage{listings}
\usepackage[most]{tcolorbox}
\usepackage{quantikz}
\usepackage[hidelinks]{hyperref}

% ---- Autor (editar a gusto) -------------------------------------------------
\newcommand{\autor}{Genesis Silesky}

% ---- Notacion de Dirac y atajos ---------------------------------------------
% quantikz ya define \ket y \bra; DeclareRobustCommand los redefine sin error.
\DeclareRobustCommand{\ket}[1]{\left|#1\right\rangle}
\DeclareRobustCommand{\bra}[1]{\left\langle#1\right|}
\DeclareRobustCommand{\braket}[2]{\left\langle#1\middle|#2\right\rangle}
\newcommand{\C}{\mathbb{C}}
\newcommand{\R}{\mathbb{R}}
\DeclareMathOperator{\Tr}{Tr}

% ---- Cajas de apoyo --------------------------------------------------------
% "concepto": definiciones previas que hacen falta para entender el capitulo.
\newtcolorbox{concepto}[1]{breakable,colback=blue!4,colframe=blue!55!black,
  fonttitle=\bfseries,title={Concepto previo: #1}}
% "nota": aclaraciones y advertencias.
\newtcolorbox{nota}[1]{breakable,colback=orange!6,colframe=orange!70!black,
  fonttitle=\bfseries,title={#1}}

% ---- Estilos de listados ----------------------------------------------------
\lstdefinestyle{py}{language=Python,basicstyle=\ttfamily\footnotesize,
  breaklines=true,postbreak=\mbox{\textcolor{gray}{$\hookrightarrow$}\space},
  numbers=left,numberstyle=\tiny\color{gray},frame=single,
  commentstyle=\color{green!40!black}\itshape,
  keywordstyle=\color{blue!70!black}\bfseries,
  stringstyle=\color{red!60!black},showstringspaces=false,tabsize=4,
  columns=fullflexible,keepspaces=true}
\lstdefinestyle{salida}{basicstyle=\ttfamily\footnotesize,breaklines=true,
  frame=leftline,backgroundcolor=\color{gray!8},columns=fullflexible,
  keepspaces=true}

  \begin{document}
  \title{\tituloconcepto}
  \author{\autor}
  \date{\today}
  \maketitle
\fi
 cada capitulo define \tituloconcepto.
\ifdefined\enmain\else
  \documentclass[11pt,a4paper]{article}
  % Preambulo compartido por todos los documentos. Compilar con XeLaTeX/LuaLaTeX
% (o Tectonic). Si se usa Overleaf, elegir el compilador XeLaTeX.
\usepackage[spanish,es-noshorthands]{babel}
\usepackage{amsmath,amssymb}
\usepackage{booktabs,array}
\usepackage[a4paper,margin=2.5cm]{geometry}
\usepackage{xcolor}
\usepackage{listings}
\usepackage[most]{tcolorbox}
\usepackage{quantikz}
\usepackage[hidelinks]{hyperref}

% ---- Autor (editar a gusto) -------------------------------------------------
\newcommand{\autor}{Genesis Silesky}

% ---- Notacion de Dirac y atajos ---------------------------------------------
% quantikz ya define \ket y \bra; DeclareRobustCommand los redefine sin error.
\DeclareRobustCommand{\ket}[1]{\left|#1\right\rangle}
\DeclareRobustCommand{\bra}[1]{\left\langle#1\right|}
\DeclareRobustCommand{\braket}[2]{\left\langle#1\middle|#2\right\rangle}
\newcommand{\C}{\mathbb{C}}
\newcommand{\R}{\mathbb{R}}
\DeclareMathOperator{\Tr}{Tr}

% ---- Cajas de apoyo --------------------------------------------------------
% "concepto": definiciones previas que hacen falta para entender el capitulo.
\newtcolorbox{concepto}[1]{breakable,colback=blue!4,colframe=blue!55!black,
  fonttitle=\bfseries,title={Concepto previo: #1}}
% "nota": aclaraciones y advertencias.
\newtcolorbox{nota}[1]{breakable,colback=orange!6,colframe=orange!70!black,
  fonttitle=\bfseries,title={#1}}

% ---- Estilos de listados ----------------------------------------------------
\lstdefinestyle{py}{language=Python,basicstyle=\ttfamily\footnotesize,
  breaklines=true,postbreak=\mbox{\textcolor{gray}{$\hookrightarrow$}\space},
  numbers=left,numberstyle=\tiny\color{gray},frame=single,
  commentstyle=\color{green!40!black}\itshape,
  keywordstyle=\color{blue!70!black}\bfseries,
  stringstyle=\color{red!60!black},showstringspaces=false,tabsize=4,
  columns=fullflexible,keepspaces=true}
\lstdefinestyle{salida}{basicstyle=\ttfamily\footnotesize,breaklines=true,
  frame=leftline,backgroundcolor=\color{gray!8},columns=fullflexible,
  keepspaces=true}

  \begin{document}
  \title{\tituloconcepto}
  \author{\autor}
  \date{\today}
  \maketitle
\fi


\section{Oscilación de Rabi: cómo un pulso cambia el estado de un qubit}

\subsection{Objetivo}
Ver cómo un pulso de radiofrecuencia resonante hace girar el estado de un qubit
de espín, por qué $P(1)$ oscila con la duración del pulso y cómo se mide la
frecuencia de Rabi $\omega_1$ en el laboratorio.

\subsection{Conceptos previos}
\begin{concepto}{espín, esfera de Bloch y evolución unitaria}
Ver capítulo~10 (esfera de Bloch) y la sección~\ref{sec:prerreq}. Idea
clave: un qubit de espín $1/2$ en un campo magnético fijo $B_0$ (eje $z$) tiene
dos niveles de energía, $\ket0$ (espín arriba) y $\ket1$ (espín abajo).
\end{concepto}

\subsection{Fundamento teórico}
Cuando la frecuencia de un campo de radiofrecuencia (RF) coincide con la del
espín (\emph{resonancia}), el espín absorbe energía del campo y cambia de
nivel. Esa es la oscilación de Rabi, y es la base del control de qubits en
estos sistemas.

\paragraph{Campo de RF y sistema rotante.} La bobina genera un campo oscilante
$B_{\mathrm{rf}}(t)\propto\cos(2\pi\omega_{\mathrm{rf}}t+\varphi_{\mathrm{rf}})$
a lo largo de $x$. Un campo linealmente polarizado se puede descomponer en dos
campos circulares, de intensidad $B_1$ cada uno, que giran en sentidos opuestos; con la \emph{aproximación de onda rotante} solo se
conserva el que gira en el mismo sentido que el espín. Si el observador gira
junto con ese campo (\emph{sistema rotante}), el campo se vuelve estático: tiene
intensidad $B_1$ y dirección
$\cos\varphi_{\mathrm{rf}}\,\hat x+\sin\varphi_{\mathrm{rf}}\,\hat y$, que
depende solo de la fase del pulso.

\paragraph{Resonancia.} En el sistema rotante, el espín gira alrededor del eje
$z$ con frecuencia $\omega_0-\omega_{\mathrm{rf}}$ (la \emph{desintonia}
$\Delta\omega$). Si $\omega_{\mathrm{rf}}=\omega_0$ ese giro desaparece y solo
queda el giro alrededor del campo $B_1$, a la \textbf{frecuencia de Rabi}
$\omega_1\propto B_1$. El ángulo de giro es
\begin{equation}
  \theta=2\pi\,\omega_1\,t_{\mathrm{Rabi}} .
\end{equation}

\paragraph{Qué se mide.} Partiendo de $\ket0$, el espín gira alrededor de un
eje del plano $xy$, así que oscila entre $\ket0$ y $\ket1$. Como $\theta$ es
también el ángulo de la esfera de Bloch,
\begin{equation}
  P(1)=\sin^2\frac\theta2=\sin^2(\pi\,\omega_1\,t),
\end{equation}
y esa población oscila con frecuencia $\omega_1$. Un pulso con $\theta=\pi$ voltea
el qubit (como la puerta $X$) y uno con $\theta=\pi/2$ lo deja en una
superposición $50/50$. Antes de usar el sistema para calcular, hay que medir
$\omega_1$: es un parámetro del equipo.

\paragraph{Cómo se mide $\omega_1$.} Medir colapsa el estado, así que:
\begin{enumerate}
  \item Se prepara siempre el estado inicial $\ket0$.
  \item Se aplica un pulso de duración $t$ y se mide muchas veces; el
    promedio da $P(1)$.
  \item Se repite aumentando poco a poco la duración $t$ del pulso, siempre
    reiniciando.
  \item Se ajusta una función seno a los datos; de ahí sale $\omega_1$.
\end{enumerate}

\paragraph{Desintonia.} Si el pulso no es resonante, el espín gira a la vez
alrededor de $z$ (a frecuencia $\Delta\omega$) y alrededor del eje del pulso.
Con una desintonia muy grande el giro alrededor de $z$ cancela el efecto del
pulso, que ya no excita al qubit.

\paragraph{Modelo usado en el código.} Con $\hbar=1$, en el sistema rotante,
\begin{equation}
  H=\frac{2\pi}{2}\Big[\Delta\omega\,\sigma_z+\omega_1\big(\cos\varphi\,\sigma_x
  +\sin\varphi\,\sigma_y\big)\Big],\qquad U(t)=e^{-iHt}.
\end{equation}
Las frecuencias están en Hz y por eso aparece $2\pi$. El capítulo~12 explica
de dónde sale $U(t)=e^{-iHt}$.

\subsection{Código (completamente comentado)}
\lstinputlisting[style=py]{../../fuentes/11_oscilacion_rabi.py}

\subsection{Explicación celda por celda}
\begin{description}
  \item[Celda 1] Importa \texttt{expm} (exponencial de una matriz),
    \texttt{curve\_fit} (ajuste de curvas) y define las matrices de Pauli.
  \item[Celda 2] Define $H$ y \texttt{estado\_tras\_pulso}, que calcula
    $U(t)\ket0$.
  \item[Celda 3] Con pulso resonante compara $P(1)$ simulada con
    $\sin^2(\pi\omega_1t)$ para cuatro duraciones.
  \item[Celda 4] Muestra que la fase del pulso elige el eje: con
    $\varphi=0$ y $\varphi=\pi/2$ un pulso $\pi/2$ deja al qubit en puntos
    distintos del ecuador.
  \item[Celda 5] Simula el experimento del laboratorio: 41 anchos de pulso,
    500 mediciones cada uno (con semilla fija), y ajusta $\omega_1$.
  \item[Celda 6] Calcula el máximo de $P(1)$ con distintas desintonías.
  \item[Celda 7] Grafica datos, teoría y el caso desafinado.
\end{description}

\subsection{Salida obtenida}
\lstinputlisting[style=salida]{../../fuentes/salidas/11_oscilacion_rabi.txt}
La simulación usa semilla fija, así que el ajuste de $\omega_1$ se repite
igual; con otra semilla cambiaría un poco.

\subsection{Interpretación}
Con $\omega_1=1000$~Hz, un pulso de $0.5$~ms da $P(1)=1$ (pulso $\pi$) y uno de
$0.25$~ms da $0.5$ (pulso $\pi/2$). El ajuste con 500 mediciones por punto
recupera $\omega_1\approx1002$~Hz. Con desintonía de $500$~Hz el máximo de
$P(1)$ baja a $0.8$, y con $5000$~Hz baja a $0.038$: el pulso casi no excita.

\subsection{Ejercicios}
\begin{enumerate}
  \item Cambia $\omega_1$ a $2000$~Hz: ¿cuánto dura ahora un pulso $\pi$?
  \item Repite el ajuste con $50$ mediciones por punto y compara el error.
  \item Con $\varphi=\pi/2$, ¿qué estado queda tras un pulso $\pi/2$?
    Compruébalo y ubícalo en la esfera de Bloch.
\end{enumerate}

\subsection{Errores comunes}
\begin{itemize}
  \item Olvidar el $2\pi$ al pasar de Hz a radianes por segundo.
  \item No reiniciar el estado entre mediciones: cada medición colapsa el
    qubit, así que el experimento se rehace desde $\ket0$.
\end{itemize}

\ifdefined\enmain\else
  \end{document}
\fi

